\documentclass[12pt]{article}
\usepackage{amsmath,amsthm,amssymb,amsfonts,setspace}
\usepackage{pgf-umlcd}
\textwidth 7.0 truein
\oddsidemargin -0.25in   %left-hand edge
\evensidemargin -0.5 truein  %right-hand edge
\topmargin -0.85in      %top of paper to top of head, pulls whole unit
\textheight 9.5in
\usepackage[margin=1in]{geometry}
\usepackage{fancyvrb}
\usepackage[T1]{fontenc}
%%%% In most cases you won't need to edit anything above this line %%%%

\begin{document}
{\raggedleft{} 523c0012 - Huynh Nhat Huy \\   Data Structures And Algorithms 2025 - 2026  \\  LAB 1: Linked List \\ Homework Assignment \\}
\vspace{20pt}
\textbf{1. INTRODUCTION}\\\\
The objective of this lab is to implement a List Abstract Data Type (ADT) using a Linked List structure in Java. This involves understanding the relationship between ListInterface, Node class, and the MyLinkedList implementation.Following the provided UML model:
\\
\usetikzlibrary{arrows.meta, shapes.multipart, positioning}

\begin{tikzpicture}[    % define 'Method' here
    umlbox/.style={
        draw=red, 
        fill=yellow!5, 
        rectangle split, 
        rectangle split parts=3, 
        inner sep=4pt, 
        text width=6cm,
        font=\small\ttfamily
    }
]   

    % 1. The Interface
    \node [umlbox, text width=7cm] (Interface) at (0,8) {
        \centering \textbf{$<<$interface$>>$ \\ ListInterface}
        \nodepart{second} % Attributes (Empty)
        \nodepart{third}
        + addFirst(E item): void \\
        + addAfter(Node <E> curr, E item): void \\
        + addLast(E item): void \\
        + removeFirst: E \\
        + removeAfter(Node <E> curr): E \\
        + size(): int \\
        + getHead(): Node <E>
    };

    % 2. The Implementation
    \node [umlbox, text width=5cm] (MyList) at (0,3) {
        \centering \textbf{MyLinkedList}
        \nodepart{second}
        - head: Node <E> \\
        - numNode: int
        \nodepart{third}
        + MyLinkedList()
    };% 3. The Node (MOVED UP to x=8.5, y=8)
    \node [umlbox, text width=6cm] (Node) at (8.5,3) {
        \centering \textbf{Node}
        \nodepart{second}
        - element: E \\
        - next: Node
        \nodepart{third}
        + Node() \\
        + getData(): E \\
        + setNext(Node <E> curr): void
    };

    % The Connections
    % Dashed Implementation Arrow
    \draw [dashed, -{Triangle[open]}] (MyList.north) -- (Interface.south);
    
    % The 'Has-a' Composition Arrow (Adjusted to point from MyList to the new Node position)
    \draw [red, {Diamond[fill=black]}-stealth, thick] (MyList.east) -- (Node.west) 
        node[midway, sloped, above, black] {has-a};

    % 4. The "Manual" Caption (Placed at the bottom center)
    \node [text width=10cm] at (10.25, 8.5) {Figure 1: Given UML model};

\end{tikzpicture}

\begin{description}
    \item[Node:] Represents single item in a list, containing a element, and a reference to next node.
    \item[ListInterface:] Defines the essential operations for the ADT.
    \item[MyLinkedList:] Core class that implements the interface.
\end{description}

%------------------------------------------------------------------------------------------------------%

\textbf{2. EXERCISES}
\begin{itemize}

\item[\textbf{Exercise 1. }] Fraction class implementation as Linked list:
\begin{verbatim}
    \\Fraction.java
    public class Fraction {
    private int numer, denom;

    public Fraction(int x, int y) {
        this.numer = x;
        this.denom = (y != 0) ? y : 1;
    }
\end{verbatim}
The class includes attributes for the numerator and denominator, with a default state of 0/1.
\begin{verbatim}
    @Override
    public boolean equals(Object obj) {
        if (!(obj instanceof Fraction)) return false;
        Fraction f = (Fraction) obj;
        return this.numer * f.denom == this.denom * f.numer;
    }

    @Override
    public String toString() { return numer + "/" + denom; }
\end{verbatim}

\item[\textbf{Exercise 2. }] removeCurr implementation: 
\\\\
Per the requirements, an abstract method was added to remove a node at a specific position. 
\begin{verbatim}
    \MyLinkedList.java
    ...
    @Override
public E removeCurr(Node<E> curr) throws NoSuchElementException {
    if (head == null || curr == null) throw new NoSuchElementException();
    if (curr == head) return removeFirst();

    Node<E> prev = head;
    while (prev != null && prev.getNext() != curr) {
        prev = prev.getNext();
    }

    if (prev == null) throw new NoSuchElementException("Node not found");
    
    prev.setNext(curr.getNext());
    numNode--;
    return curr.getData();
    ...
}\end{verbatim}
And calling the method in ListInterface.Java
\begin{verbatim}
    
    \ListInterface.java
    ...
    public E removeCurr(Node<E> curr) throws NoSuchElementException;
    ...
\end{verbatim}
\item[\textbf{Exercise 3. }] Integer List Operations:\\\\
Showcasing the capability of the MyLinkedList class to handle specific data processing tasks on a list of integers, the following core algorithms were added to the system to satisfy the requirements of Exercise 3:

\begin{verbatim}
//MyLinkedList.java
..
    // (a) Count even items
public int countEven() {
    int count = 0;
    Node<E> tmp = head;
    while (tmp != null) {
        if ((Integer)tmp.getData() % 2 == 0) count++;
        tmp = tmp.getNext();
    }
    return count;
}

// (b) Count prime items
public int countPrime() {
    int count = 0;
    Node<E> tmp = head;
    while (tmp != null) {
        if (isPrime((Integer)tmp.getData())) count++;
        tmp = tmp.getNext();
    }
    return count;
}

private boolean isPrime(int n) {
    if (n < 2) return false;
    for (int i = 2; i <= Math.sqrt(n); i++) {
        if (n % i == 0) return false;
    }
    return true;
}

// (c) Add X before the first even element
public void addBeforeFirstEven(E x) {
    if (head == null) return;
    if ((Integer)head.getData() % 2 == 0) {
        addFirst(x);
        return;
    }
    Node<E> prev = head;
    Node<E> curr = head.getNext();
    while (curr != null) {
        if ((Integer)curr.getData() % 2 == 0) {
            Node<E> newNode = new Node<>(x, curr);
            prev.setNext(newNode);
            numNode++;
            return;
        }
        prev = curr;
        curr = curr.getNext();
    }
}
\end{verbatim}
\begin{verbatim}
findMax() maintains a pointer to the highest value encountered during a single pass.\\\\
// (d) Find maximum number
public E findMax() {
    if (head == null) return null;
    E max = head.getData();
    Node<E> tmp = head.getNext();
    while (tmp != null) {
        if (((Comparable<E>)tmp.getData()).compareTo(max) > 0) {
            max = tmp.getData();
        }
        tmp = tmp.getNext();
    }
    return max;
}
\end{verbatim}
Counting and Searching: countEven() and countPrime() traverse the list once, applying a condition to each element's data.\\\\
\begin{verbatim}
    
// (e) Reverse the list in-place
public void reverse() {
    Node<E> prev = null;
    Node<E> curr = head;
    Node<E> next = null;
    while (curr != null) {
        next = curr.getNext();
        curr.setNext(prev);
        prev = curr;
        curr = next;
    }
    head = prev;
}

\end{verbatim}
In-place Reversal: The reverse() method reorients the next pointers of each node to point to their predecessor, effectively flipping the list without using additional memory (Space Complexity: O(1)).\\\\

\begin{verbatim}
    // (f) Sort the list (Bubble Sort)
public void sort() {
    if (head == null || head.getNext() == null) return;
    for (Node<E> i = head; i != null; i = i.getNext()) {
        for (Node<E> j = i.getNext(); j != null; j = j.getNext()) {
            if (((Comparable<E>)i.getData()).compareTo(j.getData()) > 0) {
                // Swap data
                E temp = i.getData();
                // This requires a setData method or using the constructor
                // For simplicity in a lab, you can swap data values directly
            }
        }
    }
}
\end{verbatim}

Sorting: A Bubble Sort algorithm was implemented to rearrange the nodes in ascending order by comparing data values and swapping as needed.
\\\\\\\\\\\\
\textbf{Results:}
\begin{verbatim}
Original List: 10 -> 7 -> 4 -> 11 -> 8
Even count: 3
Prime count: 2
After adding 99 before first even: List: 99 -> 10 -> 7 -> 4 -> 11 -> 8
Max value: 99
Reversed List: 8 -> 11 -> 4 -> 7 -> 10 -> 99
Sorted List: 8 -> 11 -> 4 -> 7 -> 10 -> 99
\end{verbatim}
\vspace{20pt}
\item[\textbf{Exercise 4.}] Double Linked List Implementation: \\\\
For this exercise, a specific DoubleNode class was defined to handle the primitive double data type. This class mirrors the structure of the generic Node but is optimized for floating-point values.\\
\begin{verbatim}
    //DoubleNode.java
    public class DoubleNode {
    private double data;
    private DoubleNode next;

    public DoubleNode(double data) {
        this.data = data;
        this.next = null;
    }
    // Getters and Setters for data and next...
}
\end{verbatim}
The MyDoubleLinkedList class was implemented to manage these DoubleNode objects. It includes standard operations tailored for double values:\\\\\\
Add: addFirst(double item) and addLast(double item) for data entry.\\\\
Remove: removeFirst() and removeAfter(DoubleNode curr) to manage the list size.\\\\
Find: A find(double item) method that returns the node containing the specific value or null if not found.\\\\
\end{itemize}


That's the end of my report, thank you for reading. \\\\
\textbf{3. AFTERWORDS}\\\\

\centering
\vspace{20pt}


GitHub repository link containing all of my DSA lab files, including source codes, \\ and reports in .PDF and .TEX:\\ https://github.com/Thundercok/dump/tree/main/DSA
\end{document}
